\section{Introduction}\label{sec:introduction}

Let $p$ be a prime, let $K(n)$ denote Morava $K$-theory of height
$n\geq 1$ at $p$, and let $S_p^\wedge$ be the $p$-completed sphere.
Morava localization isolates a single chromatic height in stable
homotopy theory. Even for the sphere, the localized spectrum is not connective, and
a degreewise finiteness statement
must control each integer degree without a connectivity assumption.
The degreewise finite-generation question asks whether each of these
groups is a finitely generated $\Zp$-module. We prove
that it is, at every prime and positive height.

\begin{theorem}\label{thm:main}\label{thm:descent-integral}
For every prime $p$, every integer $n\geq 1$, and every integer $t$,
the $\Zp$-module
\[
                 \pi_t L_{K(n)}S_p^\wedge
\]
is finitely generated. Equivalently, for every finite $p$-local
spectrum $F$, the module $\pi_t L_{K(n)}F$ is finitely generated
over $\Zp$ in every integer degree $t$.
\end{theorem}

The equivalence in the statement follows from finite cofiber sequences,
suspensions, and retracts: finite generation over the Noetherian ring
$\Zp$ is preserved by kernels, cokernels, and extensions. The proof
of the sphere assertion is completed in \cref{sec:sphere}.
There is no uniform bound asserted on the number of generators or on
the order of the torsion as $p$, $n$, or $t$ varies.

\subsection{History and relation to earlier work}
\label{sec:history}

\Cref{thm:main} answers positively the degreewise finite-generation
question stated by Hovey--Strickland \citep[Problem~16.2]{hoveyStrickland1999}
and in Hovey's Morava $K$- and $E$-theory problem list
\citep[Problem~1]{hoveyProblems}.
In Li's terminology, the question asks whether the local sphere is of
\emph{finite type}: its homotopy groups are finitely generated over $\Zp$
in every integer degree \citep[Introduction]{liFiniteType2021}.

The classical height-one calculation expresses the local sphere as the
fiber of an Adams operation on completed $K$-theory
\citep[\S\S1--2]{hopkinsKOne}; we recall the all-degree answer at odd
primes in \cref{prop:sphere-height-one}.
At height two and $p\geq5$, Hovey--Strickland's analysis of Shimomura's
calculations proves that $\pi_aL_{K(2)}F$ is finite for every finite
torsion spectrum $F$ and every integer $a$
\citep[\S15.2 and Theorem~15.1]{hoveyStrickland1999}.
Applied to the Moore spectrum, this supplies the mod-$p$ finiteness
criterion for the sphere in that range. The proof here establishes
that criterion uniformly in the prime and height.

Devinatz approached integral finiteness through continuous homotopy
fixed points for smaller subgroups of the Morava stabilizer. He used
\emph{essentially finite rank}, a condition of finite generation over
the periodicity ring after quotienting by the closure of periodicity
torsion \citep[\S4]{devinatzHomotopy2007}. For Morava $E$-theory $E_n$
and certain closed procyclic stabilizer subgroups $H$, he proved
this property for $\pi_*(E_n^{hH}\wedge F)$ under chromatic-acyclicity
and annihilation hypotheses on the finite spectrum $F$
\citep{devinatzFiniteness2008}.
Such intermediate fixed-point spectra require their own finiteness
condition; these results are distinct from degreewise finite
generation for the full local sphere.

Barthel--Schlank--Stapleton--Weinstein determined the rational answer
at every prime and height \citep[Theorem~A]{bssw2024}:
\[
 \Q\otimes_{\Z}\pi_*L_{K(n)}S
 \cong \Lambda_{\Qp}(\zeta_1,\ldots,\zeta_n),
 \qquad |\zeta_i|=1-2i.
\]
Their coefficient comparison also gives a split injection from
Witt-vector cohomology into degree-zero Lubin--Tate coefficient
cohomology whose complement in each cohomological degree is killed
by a power of $p$
\citep[Theorem~B]{bssw2024}.
These conclusions do not bound the size of integral torsion: the compact
module $\prod_{j\geq1}\Fp$ has zero rationalization and exponent $p$,
but is not finitely generated over $\Zp$.
The present theorem proves that the torsion subgroup is finite in
each integer homotopy degree.
Combined with the rational calculation, it determines the free ranks
in \cref{cor:sphere-ranks}.

Continuous Morava descent links coefficient cohomology to sphere
homotopy \citep[Theorem~1 and Appendix~A]{devinatzHopkins2004}.
The Hopkins--Ravenel smash product theorem, in the descendability
formulation of Mathew, gives a horizontal vanishing line at a finite
page \citep[Corollary~4.4, Theorem~4.18 and Proposition~10.10]{mathew2016}.
Heard states the resulting strongly convergent local Adams spectral
sequence at every prime and height, with a finite-page bound
independent of the input spectrum
\citep[Theorem~4.6 and Proposition~4.13]{heard2023}.
Consequently, finite coefficient cohomology for the Moore spectrum
is enough to obtain finite homotopy in each degree. Establishing that
coefficient finiteness is the main task of this paper.

Our geometric argument places Lubin--Tate deformation theory and
Strauch's torsion strata \citep{lubinTate1966,strauch2007} in the
geometry of the Fargues--Fontaine curve \citep{farguesFontaine2018}.
Scholze--Weinstein relate $p$-divisible groups and their universal
covers to sections of associated bundles \citep{scholzeWeinstein2013}; the relative bundle
and cohomology results used here are those of Fargues--Scholze
\citep{farguesScholze2021}.
The Banach--Colmez theory initiated by Colmez \citep{colmez2002}
was related to coherent sheaves on the curve by Le Bras
\citep{lebras2018}. Ansch\"utz--Le~Bras's relative full-faithfulness
theorem supplies the passage from the compatible linear morphisms
of section sheaves constructed below to bundle maps in families
\citep[Corollary~3.11]{anschutzLeBras2025}.
Within this framework, \cref{sec:completed-tower,sec:kummer,sec:independent-tuples}
construct the marked quotient on each height stratum and its
compact-support comparisons, retaining the marking actions and the
directions of the support limits.

The passage from geometric supports back to compact coefficients
also uses analytic-group duality. Lazard's theory and the formulation
of Beaudry--Goerss--Hopkins--Stojanoska provide the uniform-subgroup
resolutions and their orientation data \citep{lazard1965,bghs}.
The tensor-induction method, recalled by Symonds in its profinite
form, passes from an open subgroup to a complex with finite isotropy
\citep[Theorem~5.2]{symonds2007}.
In \cref{sec:continuous}, trace-one elements for these finite
stabilizers replace averaging by their orders. This permits the
complete-step coefficient comparisons and full-group norm argument
even when the stabilizer has $p$-torsion.

The transfer argument combines classical constructions with a
uniformity problem at finite marking level. Cartier's $p$-basis
construction identifies differential forms with Frobenius-twisted
de Rham cohomology \citep[\S3, Theorem~1]{cartier1957}; we use the
inverse-linear operator in the conventions of Blickle--B\"ockle
\citep[Definition~2.1 and Example~2.23]{blickleBockle2011}.
Lipman's formal residue transformation law gives the finite-root
comparison \citep[Lemmas~(7.2)(b) and~(7.3)]{lipmanResidues1984}.
The additional work is to control the poles of all transfer iterates
and to choose one finite marking level for all sufficiently high
transfers. These are the steps that return the geometric finiteness
statement to the compact coefficient modules
(\cref{sec:transfers,sec:laurent,sec:cutoff,sec:pro-transfer}).

Degreewise finite generation does not extend to every invertible
$K(n)$-local spectrum.
At height two and $p\geq5$, Hovey--Strickland's analysis of
$p$-adically graded spheres gives invertible spectra that are not of
finite type \citep[Theorem~15.1]{hoveyStrickland1999}; Li classifies
the finite-type invertibles in this range
\citep[Theorem~4.22]{liFiniteType2021}.
Thus arbitrary invertible or dualizable local spectra cannot replace
the finite spectra in \cref{thm:main}.
The proof makes no chromatic splitting assumption, and it does not
compute the finite torsion summands or the chromatic fracture maps.


\subsection{The coefficient theorem}

Fix a height $n$ formal group over a sufficiently large finite field
$\kk$ of characteristic $p$. Let $E_n$ be its Morava $E$-theory and
let $P_n$ be its ordinary stabilizer group. The algebraic conclusion
that drives the proof is
\begin{equation}\label{eq:intro-coefficients}
 H^a_{\cts}\bigl(P_n;(E_n)_{2t}/p\bigr)
 \quad\text{is finite for every }a\geq0\text{ and }t\in\Z.
\end{equation}
Here the coefficients have their compact deformation-parameter
topology. The proof retains a precise larger class: integer powers of
the periodicity line tensored with functions on whole finite tuples of
division points, with specified Frobenius-powered coordinates, and
finite sums, tensor products, equivariant summands, and finite
filtrations of these modules. Its exact definition is
\cref{def:mark-coefficients}; its finiteness theorem and Morava
consequence are \cref{thm:coeff-finite,cor:coeff-morava}.
Using whole division-point schemes allows restriction to a deeper
height stratum and descent through finite marking covers. The assertion
is about this class, rather than arbitrary semilinear finite free
modules.

Here is why this coefficient theorem suffices for the sphere.
The finite Galois quotient from $P_n$ to the extended stabilizer
$\mathbb G_n$ preserves degreewise finiteness, including when its
order is divisible by $p$. Morava descent for the Moore spectrum
$S/p$ then has finite second-page entries. A horizontal vanishing
line on a finite page leaves only finitely many of them in each
limiting stem, so $\pi_*L_{K(n)}(S/p)$ is finite in each degree.
Finally, positive-height locality gives derived $p$-completeness.
The finite $p$-power cofibers present each integral homotopy group
as a compact $\Zp$-module with finite quotient modulo $p$; compact
Nakayama proves finite generation. These steps are proved in
\cref{sec:sphere}. The horizontal bound concerns a later page of
the descent spectral sequence; finite mod-$p$ cohomological
dimension of the full stabilizer is not needed.

\subsection{The open-stratum argument}

Write
\[
 A_m=\kk[[u_m,\ldots,u_{n-1}]],\qquad
 B_m=A_m[u_m^{-1}]\quad(1\leq m<n),\qquad A_n=\kk.
\]
We prove coefficient finiteness by downward induction from $A_n$.
At the step indexed by $m$, put $x=u_m$ and $G=P_n$.
The next-height hypothesis makes the cohomology of every finite
boundary strip $x^aM/x^bM$ finite. Thus localization from a compact
allowed coefficient $M$ to $M[1/x]$ has countable kernel and
cokernel on cohomology. It remains to prove that open-stratum
cohomology is countable: the original coefficient cohomology will
then be compact Hausdorff and countable, hence finite. This explains why an
intermediate countability theorem can prove a finiteness theorem.

We work first on finite covers $B_U$ of $B_m$ that mark the
connected formal group and the \'etale Tate module. The index $U$
is an open subgroup of the two marking groups. The finite cover
is an idempotent summand of a localized whole finite free model
$A'[1/x]$. Keeping the whole model before localization ensures
that its boundary layers remain in the next-height coefficient
class. For a suitable compact lattice $\mathcal P=x^{-D}A'$, set
\[
 M^i_U=H^i_\cts(G,\mathcal P),\qquad
 X^i_U=H^i_\cts(G,B_U).
\]
The central task is countability of $X^i_U$ at every sufficiently
deep \emph{fixed} marking level. A statement only after taking the
union of all marking levels would not suffice for the return to
compact coefficients.

The completed marking tower provides the geometric input.
Its quotient by $G$ is the space of independent $r$-tuples in
$H^1(\cO(-1/m))$ on the relative Fargues--Fontaine curve, where
$\cO(-1/m)$ has rank $m$ and degree $-1$, and $r=n-m$. Uniform Kummer calculations and deletion of dependent
tuples make the equivariant compact-support cohomology finite.
A separate ordinary-function calculation shows that invariant
algebraic functions on the tower are constants. This permits a
single normalization of a Cartier transfer $S$ on $B_U$ at one
finite marking stage. Its pole estimate divides the pole order
by a fixed power of $p$, up to a fixed additive loss; consequently
$\mathcal P$ can be chosen $S$-stable and eventually contains the
transfer iterates of every bounded-pole lattice.

To use the compact-support result, we express it as a Laurent
module $W$ of convergent series with fractional $p$-power
exponents, for which every $H_i(G,W)$ is finite over $\kk$.
Residues identify
its quotient by the subspace $W_+$ of positive-$x$-order series
with the discrete transpose transfer tower:
\[
 W/W_+\simeq
 \colim\bigl(\mathcal P^\vee\xrightarrow{S^\vee}
                   \mathcal P^\vee\xrightarrow{S^\vee}\cdots\bigr).
\]
Here $\mathcal P^\vee$ is the continuous $\kk$-linear dual.
The positive-order kernel is homologically acyclic. Its homology
is countable; an analytic-subgroup argument using the Baire
property then makes boundaries open in each complete cycle space, and
absolute Frobenius contracts each positive-order cycle into that
open subgroup. Invertibility of Frobenius then makes the original
cycle a boundary. Compact duality gives
\[
 I^i_U=\bigcap_{a\geq0}S^aM^i_U\quad\text{finite}.
\]
Thus geometry controls a stable transfer image. It has not yet
controlled all of $X^i_U$.

The remaining argument uses the finite translation representations
of the generator-lift torsor. For $m>1$ their distribution algebra
is the regular local ring
$\Lambda=\kk[[D_1,\ldots,D_d]]$, with $d=(m-1)r$.
Exact pro-completion first shows that consecutive transfer images
stabilize modulo countable spaces on a marking tail. Together
with the next-height boundary comparison, this upgrades the
compact result to finiteness of
\[
 F^i_U=\bigcap_{a\geq0}
       \im\bigl(S^aM^i_U\longrightarrow X^i_U\bigr).
\]
Suppose some $X^i_U$ remained uncountable at arbitrarily deep
levels, and choose the highest such degree $b$; the finite-isotropy
bound makes this possible. A top Koszul extension then gives a
finite-diagram operation from degree $b-d$ on a fixed root cover
to degree $b$ on $B_U$, with countable cokernel. One common marking
shrink makes that operation factor through every sufficiently
high transfer. These factorizations use coinduced diagrams of
one fixed self-extension, whose boundary operator has bounded
pole loss. Modulo a countable quotient, the source classes of
the top operation have positive-order cocycle representatives.
Powering these representatives eventually absorbs that pole
loss, so their images belong to every late compact
transfer image and hence to $F^b_U$. The operation therefore has
countable image as well as countable cokernel, contradicting the
choice of $b$. For $m=1$, a Cartier splitting of Frobenius
replaces the top extension and gives the same conclusion directly.

This is the bridge from geometric finiteness to ordinary
coefficient cohomology. \Cref{sec:finite-markings,sec:completed-tower}
construct the finite models and identify the tower;
\cref{sec:kummer,sec:independent-tuples} prove its two geometric
outputs. The transfer and Laurent models are developed in
\cref{sec:transfers,sec:laurent}, and
\cref{sec:cutoff,sec:pro-transfer} prove the stable-image and
fixed-level countability assertions just described.
\Cref{sec:coefficients} extends countability to the allowed
coefficient class, descends the finite marking cover, and closes
the height induction. The continuous-cochain framework used
throughout is established first in \cref{sec:continuous}.

The geometric completion enters through compact supports and the
transpose transfer tower. Compact supports, neighborhood germs,
and supports shrinking to a closed locus use different directions
of passage; \cref{fig:support-directions} displays their maps and
the cohomology that survives. These distinctions permit the
argument to return to an individual algebraic marking level.

\subsection{Conventions}

Throughout the coefficient argument, $\kk$ is finite. It may be
enlarged by a finite extension to define the Honda models and their
endomorphisms. We write $G=P_n$ and reserve $\mathbb G_n$ for the
extended stabilizer in \cref{sec:sphere}. All group cohomology is
continuous, with the topological coefficient conventions of
\cref{sec:continuous}. Compact lattices carry their adic topology;
localized coefficients use continuous cochains with a bounded
denominator. Group homology is denoted by $H_i$ and cohomology by
$H^i$.

For analytic arguments, $C=\Cp$ and $\Cflat$ is its tilt. If
$\mathcal D$ is a compact profinite parameter space, write
$R_{\mathcal D}=C(\mathcal D,\Cflat)$ for its continuous-function
ring. Statements of finiteness with parameters mean finiteness
over $R_{\mathcal D}$. The coefficient-only Frobenius $\varphi$
fixes monomial exponents; the absolute Frobenius $\Fr$ also scales
them. These two operations are distinguished throughout.

We use cohomological grading for complexes. A limit over shrinking
supports, a colimit of restrictions defining a germ, and a colimit
of corestrictions over expanding supports are specified separately
whenever they occur.
